\chapter{绪论}\label{Chap:chap1}

\section{研究背景与动机}

随着云计算、大数据、人工智能、移动互联网和物联网技术的蓬勃发展，端到端网络传输的服务对象正在从固定主机扩展到手机、笔记本、台式机、边缘节点和传感设备等多模态终端，传输路径也从近端局域链路延伸到跨城、跨域、蜂窝接入、无线局域网和卫星通信等远距离复杂网络。传输控制协议（Transmission Control Protocol，TCP）作为互联网传输层的核心协议，承载了大量关键业务流量，其拥塞控制机制的性能直接影响网络资源利用效率、应用响应延迟和用户服务质量。网络拥塞控制是计算机网络领域最具挑战性的研究课题之一，自互联网诞生以来一直是学术界和工业界持续关注的焦点问题。

在现代端到端网络中，链路与终端条件呈现出更强的异构性。一方面，城域/广域网、蜂窝网络、WiFi与卫星链路具有不同的RTT量级、丢包背景、可用带宽和链路抖动；另一方面，手机、电脑、边缘设备等终端在处理能力、无线接入方式、移动性和业务负载上差异显著。这种长距离、多接入、多终端并存的网络环境对传统拥塞控制算法提出了前所未有的严峻挑战。

现代网络应用承载着多种类型的业务负载，这些业务对网络性能有着截然不同的需求。移动视频、实时会议和云游戏需要稳定的低延迟与低抖动；分布式存储、内容分发和大文件同步需要高吞吐量以支持大规模数据传输；边缘计算、工业物联网和车联网则同时关注可靠性、实时性和终端资源约束。单一的拥塞控制策略难以同时满足这些多样化的性能需求。

更为复杂的是，远距离和多接入传输中的拥塞信号更难解释。丢包可能来自真实队列溢出，也可能来自无线误码、链路切换或终端侧调度抖动；RTT升高既可能表示排队，也可能来自路径变化、接入链路波动或终端处理延迟；移动终端后台同步、视频传输、RPC突发、备份任务和多租户共享链路还会造成短时间带宽挤占。传统单一信号驱动的拥塞控制方法难以在这些场景中稳定区分拥塞程度。

传统的TCP拥塞控制算法在这种复杂的网络环境中显露出诸多不足。基于丢包的算法（如TCP NewReno\upcite{newreno}、TCP CUBIC\upcite{cubic}）将数据包丢失作为拥塞信号，但丢包是拥塞的滞后指示器，通常意味着网络缓冲区已经溢出，排队延迟已经积累到较高水平。基于时延的算法（如TCP Vegas\upcite{vegas}）和基于带宽/最小RTT模型的TCP BBR\upcite{bbr}尝试在丢包前调节发送行为，但RTT与带宽估计易受路径变化影响，与基于丢包算法共存时的公平性也取决于具体配置。

显式拥塞通知（Explicit Congestion Notification，ECN）机制为解决上述问题提供了新的思路。ECN允许网络设备在队列长度超过阈值时标记数据包而非丢弃，从而实现更早、更精确的拥塞信号传递。数据中心TCP（DCTCP\upcite{dctcp}）验证了ECN细粒度反馈在特定低时延网络中的有效性，但其适用前提较强；在远距离、多接入和多模态终端并存的环境中，现有算法对ECN、RTT、丢包和协议栈状态的联合利用仍不充分，大多采用静态参数配置，缺乏对动态网络环境的自适应能力。

近年来，机器学习特别是强化学习技术为网络拥塞控制提供了新的研究思路\upcite{rlcc}：Aurora、Orca等方案通过训练神经网络策略或在线学习调整发送行为，具有适应不同网络条件的潜力。然而，学习型方案通常需要训练与模型推理设施，其训练分布外泛化能力仍需专门验证，决策过程也较难逐事件审计。部分方案主要依赖窗口、时延和吞吐统计，未直接利用ECN与拥塞控制状态机等协议栈内部信号。

本文因此选择\emph{白盒启发式}的设计路线：不引入学习型策略，而是在协议栈回调点采集包含ECN、拥塞状态机与窗口/RTT在内的多维状态，以确定性规则在确认或拥塞事件上输出拥塞窗口与慢启动阈值决策，仅在参数层使用基线相对的轻量反馈自适应。该路线侧重可解释性和可重复性，并通过多信号融合与参数自适应处理不同路径条件。

\section{研究目标与主要贡献}

针对上述问题，本文提出Swift算法，一种面向远距离传输、多接入链路和异构终端场景的启发式多信号融合自适应网络拥塞控制方法。其设计目标是在给定安全边界内提高链路利用率，并控制排队时延和窗口振荡。

图\ref{fig:swift_contribution_overview}展示了Swift算法的三项核心贡献及各组件之间的逻辑关系。

\begin{figure}[htbp]
\centering
\begin{tikzpicture}[
    node distance=1.0cm and 1.2cm,
    block/.style={rectangle, draw, fill=blue!8, text width=4.6cm, minimum height=1.0cm, align=center, rounded corners=3pt, font=\small},
    arrow/.style={->, >=stealth, thick, draw=gray!70},
]
% 顶层目标
\node[block, fill=red!12, text width=10cm, minimum height=0.9cm] (goal) {Swift设计目标：面向远距离传输与多模态终端的高吞吐量 $+$ 低延迟 $+$ 稳定性};

% 第二层：三项核心创新
\node[block, below left=1.0cm and -0.2cm of goal] (obs) {贡献1：多信号状态感知与分类\\{\scriptsize 11维有效观测，区分超时、ECN与普通丢包}};
\node[block, below=1.0cm of goal] (detect) {贡献2：启发式反馈自适应融合控制\\{\scriptsize 时间窗口BDP估计、路径感知调参与有界目标跟踪}};
\node[block, below right=1.0cm and -0.2cm of goal] (fusion) {贡献3：稳定性与安全约束\\{\scriptsize 差异化保留、连续缩减保护、冻结与窗口边界}};

% 底层
\node[block, below=1.2cm of detect, fill=green!8, text width=10cm] (impl) {实现验证：ns-3 $+$ ZeroMQ同步信道 $+$ 多场景对照实验};

% 箭头
\draw[arrow] (goal.south) -- ++(0,-0.25) -| (obs.north);
\draw[arrow] (goal.south) -- (detect.north);
\draw[arrow] (goal.south) -- ++(0,-0.25) -| (fusion.north);
\draw[arrow] (obs.south) -- ++(0,-0.35) -| (impl.north);
\draw[arrow] (detect) -- (impl);
\draw[arrow] (fusion.south) -- ++(0,-0.35) -| (impl.north);
\end{tikzpicture}
\caption{Swift算法核心贡献概览}
\label{fig:swift_contribution_overview}
\end{figure}

本文的主要研究贡献如下：

\textbf{（1）面向远距离与异构终端的多信号拥塞状态感知}

本文设计了一个11维有效观测子空间（嵌入于15元素状态传输容器，另含套接字UUID、环境类型、仿真时间、节点ID四项跨进程路由元数据），将ECN状态、拥塞控制状态机状态、拥塞事件类型、RTT和在途字节等信息纳入状态表示，并据此区分超时、ECN驱动的窗口缩减与普通丢包。该观测子空间为不同接入与路径条件下的拥塞判定提供协议语义；本文的仿真场景用于抽象这些链路条件，不等同于具体终端硬件测试。

\textbf{（2）启发式反馈自适应的融合控制}

本文采用滑动时间窗口内的累计确认量估计交付速率，经窗口化最大值滤波并结合最小RTT得到BDP估计；控制器根据最小RTT感知的RTT膨胀阈值、性能反馈值相对慢速基线的偏移和连续增长趋势调节每流乘性增加因子$\alpha\in[0.85,1.30]$，并以有界步长跟踪$\alpha\times BDP$目标窗口。该过程由确定性规则执行，不包含策略训练或神经网络推理。

\textbf{（3）面向不确定拥塞信号的稳定性与安全约束}

Swift对普通丢包、ECN和超时分别采用0.70、0.75和0.50的窗口保留因子，并设置最多3次连续缩减、降窗后4个ACK事件冻结、窗口上下界和CA\_LOSS下的陈旧动作作废。上述机制用于限制连续过度缩减、快速反弹与陈旧决策覆盖，具体独立贡献仍需消融实验验证。

\section{论文组织结构}

本文的其余部分组织如下：

第2章介绍相关工作，包括传统拥塞控制算法、基于机器学习的拥塞控制方法、显式拥塞通知机制、智能决策方法在网络领域的应用综述等。

第3章详细描述Swift算法的设计与实现，包括系统架构、观测空间设计、多信号融合拥塞检测、差异化拥塞响应、自适应参数调优、融合控制策略等核心组件。

第4章介绍实验评估，包括实验环境配置、评测指标、对比算法、实验场景设计以及详细的实验结果分析。

第5章讨论Swift算法的适用边界和未来研究方向。

第6章总结全文并展望未来工作。
